\section{System Model and Threat Model}
\label{sec:system-model}

\subsection{Entities}

\sysname has six entities. The data subject $\mathit{DS}_i$ controls
the private key of one soulbound Consent NFT. The data controller
$\mathit{DC}$ orchestrates the FL round loop and triggers unlearning
on revocation events. Each FL participant device $\mathit{FL}_k$ runs
local SGD and submits per-round updates. The Verifier Committee
$\mathit{VC} = \{v_1, \ldots, v_n\}$ holds $n=5$ registered ECDSA
keys; threshold $t=3$ by default. Three on-chain contracts
--- \texttt{ConsentRegistry} (CR), \texttt{UnlearningOracle} (UO), and
\texttt{FLGatekeeper} (GK) --- collectively implement the protocol.
Anyone with chain read access is an \emph{auditor}; public
auditability is a first-class property.

\subsection{Federated learning baseline}

We adopt FedAvg~\cite{mcmahan2017fedavg} as the canonical aggregation
rule. Let $\theta^{(r)}$ denote the global parameters at round $r$,
$K_r \subseteq \{1,\ldots,K\}$ the set of validated participants
returned by $\mathit{GK}.\mathtt{validateRound}$, and $n_k$ the local
training-set size of client $k$:
\begin{equation}\label{eq:fedavg}
    \theta^{(r+1)} \;=\; \sum_{k \in K_r}
        \frac{n_k}{\sum_{j \in K_r} n_j} \,\theta^{(r+1)}_k
\end{equation}
The novel restriction is that $K_r$ must be the
\texttt{validateRound} return value rather than an
orchestrator-chosen subset --- consent is enforced at the
$K_r$ membership level, not at the gradient level.

\subsection{Consent Lifecycle State Machine}
\label{sec:clsm-formal}

\begin{definition}[\clsm]\label{def:clsm}
The Consent Lifecycle State Machine is the deterministic finite
automaton $\clsm = (S, \Sigma, \delta, s_0, F)$ where:
\begin{align*}
S &= \{\textsc{Inactive}, \textsc{Active}, \textsc{Suspended}, \\
  &\qquad\textsc{Revoked}, \textsc{Expired}, \textsc{Unlearned}\}, \\
\Sigma &= \{\mathtt{grant}, \mathtt{modify}, \mathtt{suspend},
    \mathtt{resume}, \mathtt{revoke}, \\
  &\qquad \mathtt{tick}, \mathtt{expire}, \mathtt{unlearn\_verified}\}, \\
s_0 &= \textsc{Inactive},\qquad F = \{\textsc{Unlearned}\}.
\end{align*}
$\delta: S \times \Sigma \to S$ is a partial function whose transition
diagram appears in Figure~\ref{fig:clsm} and whose tabular form is
Table~\ref{tab:clsm} in Section~\ref{sec:framework-clsm};
\textsc{Unlearned} is the absorbing terminal state.
\end{definition}

\begin{figure}[t]
\centering
\scalebox{0.82}{
\begin{tikzpicture}[
    > = {Stealth[length=2mm]},
    state/.style = {circle, draw, minimum size=11mm,
                     align=center, font=\footnotesize, inner sep=1pt},
    every edge/.style = {draw, ->, font=\scriptsize},
]
\node[state]                                  (I) {\textsc{In-}\\\textsc{active}};
\node[state, right=22mm of I]                 (A) {\textsc{Active}};
\node[state, above=12mm of A]                 (S) {\textsc{Susp.}};
\node[state, below=12mm of A]                 (R) {\textsc{Revo-}\\\textsc{ked}};
\node[state, right=22mm of A]                 (E) {\textsc{Exp.}};
\node[state, double, right=22mm of R]         (U) {\textsc{Un-}\\\textsc{learned}};

\path[->]
    (I) edge node[above]              {\texttt{grant}}       (A)
    (A) edge[loop left] node          {\texttt{tick}}        (A)
    (A) edge[bend left=35]  node[right=1pt] {\texttt{susp}}   (S)
    (S) edge[bend left=35]  node[left=1pt]  {\texttt{resume}} (A)
    (A) edge[bend right=35] node[right=1pt] {\texttt{revoke}} (R)
    (A) edge                node[above]     {\texttt{expire}} (E)
    (R) edge                node[below]     {\texttt{uv}}     (U)
    (E) edge                node[right, pos=0.45] {\texttt{uv}} (U);
\end{tikzpicture}}
\caption{Consent Lifecycle State Machine.
$\mathtt{modify}$ shares the \textsc{Active} self-edge with
$\mathtt{tick}$; $\mathtt{uv} \!=\! \mathtt{unlearn\_verified}$
is callable only by the oracle.
\textsc{Unlearned} (double circle) is absorbing.}
\label{fig:clsm}
\end{figure}

The on-chain contract \texttt{ConsentRegistry} is an exact
instantiation of \clsm: each consent NFT carries a \texttt{ConsentState}
enum stored in slot 0 of the \texttt{ConsentRecord} struct, every
$\delta$ edge corresponds to one external function, and the soulbound
constraint (no \texttt{transferFrom}) enforces that only $\mathit{DS}_i$
may invoke transitions on token $i$.

\subsection{Threat model}

We adopt the adversary set summarised in
Table~\ref{tab:threat}. Five in-scope adversaries are defended against
by specific protocol mechanisms; six classes of attack are
acknowledged as out-of-scope with the rationale stated in
Section~\ref{sec:discussion}. The expanded model with per-adversary
capability/intent/defense/residual-risk lives in the supplementary
\texttt{docs/threat-model.md}.

\begin{table}[t]
\centering
\caption{Threat model summary (full version in supplementary).}
\label{tab:threat}
\small
\begin{tabular}{ll}
\toprule
Adversary & Defense \\
\midrule
Lazy FL Orchestrator & GK reverts on pending unlearning; \\
                     & $\dmax$ timeout emits ComplianceViolation \\
Deceptive Controller & Three MIAs + VC $t$-of-$n$ signatures \\
Colluding participants & Out of scope (orthogonal DP layer) \\
Consent forger & Soulbound NFT + \texttt{dsOwner} check \\
Replay attacker & UA binds (chainid, tokenId, hashPre) \\
\bottomrule
\end{tabular}
\end{table}

\subsection{Formal properties}
\label{sec:properties}

We state four properties that \sysname must satisfy. Each is proved
informally in Section~\ref{sec:framework}; full mechanised proofs are
future work.

\begin{property}[Consent Integrity --- safety]\label{prop:safety}
For every block height $b$ and round $r$ that starts at $b$ via
$\mathit{GK}.\mathtt{validateRound}$,
\begin{equation}\label{eq:safety}
    \forall i \in K_r:\quad \mathit{CR}.\mathtt{state}(i) =
    \textsc{Active}_{\,b} \;\wedge\;
    \rho_r \in \mathcal{P}_i,
\end{equation}
where $\rho_r$ is the round purpose passed to
\texttt{validateRound} and $\mathcal{P}_i$ is the purpose set of the
consent token $i$ recorded at \texttt{grantConsent} time. No transition
path in $\delta$ leads from \textsc{Revoked} back to \textsc{Active}.
\end{property}

\begin{property}[Unlearning Enforceability --- liveness]\label{prop:liveness}
For every revocation event recorded at block $b$ on token $i$,
\begin{equation}\label{eq:liveness}
\begin{aligned}
    \exists\, b' \!\le\! b\!+\!\dmax:\;
    &\mathit{CR}.\mathtt{state}(i)\!=\!\textsc{Unlearned}_{b'} \\
    &\;\;\vee\;\;
    \mathit{GK}.\mathtt{ComplianceViolation}(i)_{b'}.
\end{aligned}
\end{equation}
Equivalently: either a valid UA is committed within $\dmax$ blocks of
the revocation, or the chain logs a publicly verifiable
ComplianceViolation event that any auditor can pick up. The system
degrades to ``no training'' rather than ``unconsented training.''
\end{property}

\begin{property}[Unlearning Verifiability]\label{prop:verifiability}
The \texttt{UnlearningOracle} accepts a UA submission $(u, \Sigma)$
if and only if
\begin{equation}\label{eq:verifiability}
    \big(\forall j \in \{\mathrm{sh}, \mathrm{lira}, \mathrm{lo}\}\!:\;
        s_j < \tauMIA\big) \;\wedge\;
    \big|\{\sigma \in \Sigma\!:\!\mathrm{verify}(\sigma, h_u, \mathit{VC})\}\big| \ge t,
\end{equation}
where $s_j$ are the three MIA scores in the UA, $h_u$ is the UA hash
(Eq.~\ref{eq:uahash}), and verifications are over the registered
verifier set $\mathit{VC}$ of size $n$.
\end{property}

\begin{property}[Audit Completeness]\label{prop:audit}
Every consent state transition, every Unlearning Attestation
acceptance, every ComplianceViolation, and every per-round model hash
update emits a Solidity event indexed by token ID and round ID.
\end{property}

Properties~\ref{prop:safety} and \ref{prop:liveness} are the safety
and liveness halves of the classical agreement-protocol decomposition.
Property~\ref{prop:verifiability} is the cryptographic gating
condition. Property~\ref{prop:audit} ensures that the chain serves as
the durable evidence layer required by GDPR Article~5(2) accountability
principle and EU AI Act Article~12 record-keeping. The empirical
evidence for \ref{prop:liveness} and \ref{prop:verifiability} is
reported in Section~\ref{sec:evaluation}; the on-chain enforcement of
\ref{prop:safety} and \ref{prop:audit} follows by construction from
the contract code (Section~\ref{sec:framework}).
